\section{Nil ideals with uniquely divisible additive groups}

There are positive mixed-characteristic examples outside the coherent
case.  The reason is elementary rigidity: finite-group cohomology
vanishes on a module on which division by the group order is possible.
The coefficient ring itself need not contain that inverse.

\begin{theorem}\label{thm:rational-nil}
Let $A$ be a commutative local ring and $I$ a nil ideal whose
additive group is uniquely divisible.  For every finite group $G$,
reduction induces an isomorphism
\[
 K_0(\Lat_A(G))\xrightarrow{\ \sim\ }
 K_0(\Lat_{A/I}(G)).
\]
These isomorphisms commute with tensor products, restriction and
induction.  Consequently $\Phi_A$ is an isomorphism if and only
if $\Phi_{A/I}$ is an isomorphism.
\end{theorem}

Here a nil ideal means that every element is nilpotent; a common
nilpotence exponent for the ideal is not required.  Unique divisibility
means precisely that its additive group is a $\Q$-vector space.

\begin{proof}
We first suppose $I^2=0$ and put $B=A/I$ and $N=|G|$.
Every matrix in $\mathrm{GL}_r(B)$ lifts to one in
$\mathrm{GL}_r(A)$, since an element invertible modulo a nil ideal
is invertible.  Given a representation $\bar\rho$ on $B^r$,
choose lifts $s_g\in\mathrm{GL}_r(A)$, with $s_1=1$.
The kernel of the reduction map on general linear groups is
$1+M_r(I)$; its multiplication is addition of the matrices in
$M_r(I)$.  Conjugation by $s_g$ gives an action of $G$ on
$M=M_r(I)$, independent of the chosen lifts.

Define $c(g,h)\in M$ by
\[
 s_gs_h=(1+c(g,h))s_{gh}.
\]
Associativity gives
\[
 g\,c(h,k)-c(gh,k)+c(g,hk)-c(g,h)=0.
\]
Division by $N$ makes sense in $M$.  Put
$b(g)=N^{-1}\sum_{k\in G}c(g,k)$.  Summing the last identity
over $k$ gives
\[
 c(g,h)=b(g)+g\,b(h)-b(gh).
\]
Replacing $s_g$ by $(1-b(g))s_g$ therefore gives a representation
lifting $\bar\rho$.

Two such lifts are isomorphic by an isomorphism reducing to the
identity.  Indeed write $\rho'(g)=(1+d(g))\rho(g)$.
The homomorphism identities give
$d(gh)=d(g)+g\,d(h)$.  If
$b=N^{-1}\sum_h d(h)$, summation gives $d(g)=b-gb$.
Conjugation by $1+b$ carries $\rho$ to $\rho'$.
An arbitrary isomorphism between reductions can first be lifted
as an invertible matrix and then corrected in this manner.
Thus reduction is a bijection on isomorphism classes of lattices.

The same argument lifts every exact-sequence relation.  Choose
bases over $B$ adapted to a split underlying sequence
$0\to B^a\to B^{a+b}\to B^b\to0$.
Its group action lies in the group of invertible block upper
triangular matrices.  Reduction of that group is surjective;
its kernel is the additive group of block upper triangular
matrices with entries in $I$.  This is again uniquely divisible,
and the preceding two-cocycle calculation applies unchanged.
The resulting block representation over $A$ gives an admissible
exact sequence lifting the original one.  Conversely reduction
preserves every admissible exact sequence, because it splits over
$A$.  A bijection on isomorphism classes together with lifting
of these defining relations proves the asserted $K_0$ isomorphism.
No assertion of full faithfulness on morphisms is needed.

If $I$ is nilpotent, apply the square-zero case successively to
$A/I^{j+1}\to A/I^j$.  Each quotient $I^j/I^{j+1}$ is
uniquely divisible.  To see this, division by an integer in $I$
is $A$-linear, and
$(x_1\cdots x_j)/m=(x_1/m)x_2\cdots x_j$.
Thus every $I^j$ is a $\Q$-subspace of $I$.

Finally suppose merely that $I$ is nil.  Any finite set of its
elements generates a nilpotent ideal $J$, since the ring is
commutative.  Its rational saturation
\[
 J_{\Q}=\{x\in I:mx\in J\text{ for some integer }m>0\}
\]
is an $A$-ideal, is uniquely divisible, and is nilpotent with
the same bound: a product of sufficiently many elements becomes
zero after multiplication by a positive integer, and $I$ is
torsion-free.
For a proposed lift of a finite-group representation, only
finitely many entries occur in its multiplication defects.
Put them in such a $J_{\Q}$.  The matrices then define a
representation over $A/J_{\Q}$, and the nilpotent case lifts it
to $A$.  For uniqueness, put the finitely many entries of the
difference of two lifts, after a lifted change of basis, into
such an ideal.  The same argument works in the block upper
triangular groups and lifts every exact-sequence relation.
This proves the general nil case.

Reduction commutes with all three stated operations.  Applying
the result to every symmetric group and its induction product
therefore gives the final assertion.
\end{proof}

\begin{example}\label{ex:noncoherent-positive}
For a prime $p$ let
\[
 A=\Z_{(p)}\oplus\Q\varepsilon,\qquad\varepsilon^2=0,
\]
with the usual scalar action of $\Z_{(p)}$ on $\Q$.
This is a local ring with maximal ideal
$p\Z_{(p)}\oplus\Q\varepsilon$.  Its square-zero ideal
$\Q\varepsilon$ is uniquely divisible, and its quotient is
$\Z_{(p)}$.  Hence its induction ring is polynomial on the
trivial classes.  It is not coherent: the annihilator of
$\varepsilon$ is $\Q\varepsilon$, which is not finitely
generated as an $A$-module, whereas annihilators of individual
elements are finitely generated in a coherent ring.
This example also shows why positive coefficient rings cannot
be characterized simply by coherence and regularity.
\end{example}
